\documentclass[12pt,a4paper]{article}
\usepackage[utf8]{inputenc}
\usepackage[T1]{fontenc}
\usepackage{amsmath,amssymb,amsthm}
\usepackage{graphicx}
\usepackage{xcolor}
\usepackage{hyperref}
\usepackage{geometry}
\usepackage{tcolorbox}
\usepackage{needspace}
\usepackage{tikz}
\usepackage{array}
\usepackage{booktabs}
\usepackage{pifont}
\usetikzlibrary{arrows,positioning,matrix}
\geometry{margin=2.5cm}

\newcommand*\circled[1]{\tikz[baseline=(char.base)]{\node[shape=circle,draw,inner sep=2pt] (char) {#1};}}

\definecolor{titrebleu}{RGB}{0,102,204}
\definecolor{defvert}{RGB}{0,153,76}
\definecolor{exempleorange}{RGB}{255,102,0}
\definecolor{algoviolet}{RGB}{153,0,153}
\definecolor{importantrouge}{RGB}{204,0,0}
\definecolor{fondgris}{RGB}{245,245,245}
\definecolor{fondjaune}{RGB}{255,250,205}
\definecolor{fondvert}{RGB}{230,255,230}
\definecolor{fondbleu}{RGB}{230,240,255}
\definecolor{fondrose}{RGB}{255,230,240}

\newtcolorbox{definition}{colback=fondvert,colframe=defvert,fonttitle=\bfseries,title=Definition,arc=3mm}
\newtcolorbox{exemple}{colback=fondjaune,colframe=exempleorange,fonttitle=\bfseries,title=Exemple,arc=3mm}
\newtcolorbox{important}{colback=white,colframe=importantrouge,fonttitle=\bfseries,title=Important,arc=3mm}
\newtcolorbox{astuce}{colback=fondbleu,colframe=titrebleu,fonttitle=\bfseries,title=A savoir,arc=3mm}
\newtcolorbox{methode}{colback=fondbleu,colframe=algoviolet,fonttitle=\bfseries,title=Methode,arc=3mm}
\newtcolorbox{notecours}{colback=fondrose,colframe=importantrouge,fonttitle=\bfseries,title=Note de Cours,arc=3mm}
\newtcolorbox{theoreme}{colback=fondgris,colframe=algoviolet,fonttitle=\bfseries,title=Theoreme,arc=3mm}

\usepackage{titlesec}
\titleformat{\section}{\color{titrebleu}\normalfont\Large\bfseries}{\color{titrebleu}\thesection}{1em}{}
\titleformat{\subsection}{\color{defvert}\normalfont\large\bfseries}{\color{defvert}\thesubsection}{1em}{}

\title{\textbf{\Huge\color{titrebleu}La Methode du Simplexe Revise}\\[0.5em]
\large\color{defvert}Description Matricielle, Iterations et Interpretation Economique\\
\small Modeles Lineaires de la Recherche Operationnelle}
\author{\large Prof. Herve KERIVIN, PhD\\
\small Universite Clermont Auvergne\\
Clermont Auvergne INP - ISIMA\\
LIMOS UMR 6158 CNRS}
\date{\large Version 2024-2025 - Cours 6}

\begin{document}

\maketitle
\tableofcontents
\newpage

%=============================================
\section{Motivation : Pourquoi le Simplexe Revise ?}
%=============================================

\needspace{10\baselineskip}
\subsection{Limites du Simplexe Standard}

\begin{notecours}
\textbf{Rappel -- Le simplexe standard :}

La methode du simplexe standard met a jour un \textcolor{titrebleu}{\textbf{dictionnaire}} (ou un tableau) a chaque iteration. Cela implique de recalculer et stocker l'integralite de la matrice transformee.

Pour des problemes de grande taille en pratique (des milliers de variables et de contraintes), cette approche est \textcolor{importantrouge}{\textbf{couteuseuse en memoire et en calcul}}.
\end{notecours}

\begin{astuce}
\textbf{Observation fondamentale :}

\begin{itemize}
    \item A chaque iteration, seule une \textcolor{defvert}{\textbf{petite partie}} du dictionnaire est reellement utilisee
    \item La partie utilisee peut etre \textcolor{titrebleu}{\textbf{reconstruite a partir des donnees originales}} a tout moment
    \item Il suffit de connaître la base courante $B$ !
\end{itemize}
\end{astuce}

\begin{definition}
\textbf{Methode du simplexe revise :}

Implementation de la methode du simplexe qui travaille avec les \textcolor{defvert}{\textbf{donnees originales du probleme}} (matrice $A$, vecteurs $b$ et $c$) et la base courante $B$, \textcolor{importantrouge}{\textbf{sans maintenir de dictionnaire ou tableau transforme}}.
\end{definition}

\begin{important}
\textbf{Avantages en pratique :}

\begin{itemize}
    \item Sur les problemes \textcolor{titrebleu}{\textbf{grands et creux}} (sparse) typiques des applications, le simplexe revise est \textbf{beaucoup plus rapide} que le simplexe standard
    \item \textcolor{defvert}{\textbf{Tous les logiciels modernes}} de resolution de PL utilisent une forme du simplexe revise
    \item Chaque iteration requiert la resolution de \textcolor{algoviolet}{\textbf{deux systemes lineaires}} (pas de zero) plutot que la mise a jour complete du tableau
    \item Les matrices creuses restent creuses (pas de fill-in artificiel)
\end{itemize}
\end{important}

\newpage
%=============================================
\section{Description Matricielle du Simplexe}
%=============================================

\needspace{10\baselineskip}
\subsection{Formulation Matricielle du PL}

\begin{methode}
\textbf{Mise en forme matricielle du probleme standard :}

Considerons le PL sous forme standard :
\begin{align*}
\text{maximiser} \quad & z = \sum_{j=1}^{n} c_j x_j\\
\text{sujet a} \quad & \sum_{j=1}^{n} a_{ij} x_j \leq b_i \quad \text{pour } i = 1, 2, \ldots, m\\
& x_j \geq 0 \quad \text{pour } j = 1, 2, \ldots, n
\end{align*}

Apres introduction des variables d'ecart $x_{n+1}, \ldots, x_{n+m} \geq 0$, le probleme devient :
$$\text{maximiser } z = \mathbf{c}^T \mathbf{x} \quad \text{sujet a } A\mathbf{x} = \mathbf{b}, \quad \mathbf{x} \geq \mathbf{0}$$

ou :
\begin{itemize}
    \item $A$ est une matrice $m \times (n+m)$ : les $n$ premieres colonnes sont les contraintes originales, les $m$ dernieres forment la \textcolor{defvert}{\textbf{matrice identite}} $I_m$ (colonnes des variables d'ecart)
    \item $\mathbf{x} \in \mathbb{R}^{n+m}$ : vecteur de toutes les variables
    \item $\mathbf{b} \in \mathbb{R}^m$ : membres droits
    \item $\mathbf{c} \in \mathbb{R}^{n+m}$ : coefficients objectif (les $m$ derniers sont nuls)
\end{itemize}
\end{methode}

\begin{exemple}
\textbf{Exemple de reference :}

\begin{align*}
\text{maximiser} \quad & z = 5x_1 + 4x_2 + 3x_3\\
\text{sujet a} \quad & 2x_1 + 3x_2 + x_3 \leq 5\\
& 4x_1 + x_2 + 2x_3 \leq 11\\
& 3x_1 + 4x_2 + 2x_3 \leq 8\\
& x_1, x_2, x_3 \geq 0
\end{align*}

Avec les variables d'ecart $x_4, x_5, x_6$ :

$$A = \begin{pmatrix} 2 & 3 & 1 & 1 & 0 & 0 \\ 4 & 1 & 2 & 0 & 1 & 0 \\ 3 & 4 & 2 & 0 & 0 & 1 \end{pmatrix}, \quad \mathbf{b} = \begin{pmatrix} 5 \\ 11 \\ 8 \end{pmatrix}, \quad \mathbf{c} = \begin{pmatrix} 5 \\ 4 \\ 3 \\ 0 \\ 0 \\ 0 \end{pmatrix}$$
\end{exemple}

\newpage
\subsection{Partition Base / Hors-Base}

\begin{definition}
\textbf{Partition associee a une base $B$ :}

Chaque solution de base realisable induit une partition de $\mathbf{x}$, $A$ et $\mathbf{c}$ en parties basique et non-basique :

\begin{center}
\begin{tabular}{l|l|l}
\toprule
\textbf{Element} & \textbf{Partie basique} & \textbf{Partie non-basique} \\
\midrule
Variables & $\mathbf{x}_B$ ($m$ variables) & $\mathbf{x}_N$ ($n$ variables) \\
Colonnes de $A$ & $B = A_B$ (matrice $m \times m$) & $A_N$ (matrice $m \times n$) \\
Coefficients objectif & $\mathbf{c}_B$ ($m$ composantes) & $\mathbf{c}_N$ ($n$ composantes) \\
\bottomrule
\end{tabular}
\end{center}

Ainsi, le systeme $A\mathbf{x} = \mathbf{b}$ s'ecrit :
$$B\mathbf{x}_B + A_N\mathbf{x}_N = \mathbf{b}$$
\end{definition}

\begin{definition}
\textbf{Matrice de base (Basis Matrix) :}

La \textcolor{defvert}{\textbf{matrice de base}} $B$ de $A$ est une sous-matrice carree formee de $m$ colonnes \textcolor{algoviolet}{\textbf{lineairement independantes}} de $A$.

\begin{itemize}
    \item Une matrice de base est \textcolor{titrebleu}{\textbf{non singuliere}} (inversible)
    \item La colonne de $B$ correspondant a la variable de base $x_i$ est la $i$-eme colonne de $A$
    \item Par convention, on note $B$ la matrice de base plutot que $A_B$
\end{itemize}
\end{definition}

\begin{methode}
\textbf{Expression du dictionnaire en termes matriciels :}

En multipliant $B\mathbf{x}_B = \mathbf{b} - A_N\mathbf{x}_N$ par $B^{-1}$ a gauche :

$$\boxed{\mathbf{x}_B = B^{-1}\mathbf{b} - B^{-1}A_N\mathbf{x}_N}$$

Et en substituant dans $z = \mathbf{c}^T\mathbf{x} = \mathbf{c}_B^T\mathbf{x}_B + \mathbf{c}_N^T\mathbf{x}_N$ :

$$\boxed{z = \mathbf{c}_B^T B^{-1}\mathbf{b} + (\mathbf{c}_N^T - \mathbf{c}_B^T B^{-1}A_N)\mathbf{x}_N}$$

Ces deux equations constituent le \textcolor{algoviolet}{\textbf{dictionnaire en forme matricielle}}, associe a la base $B$.
\end{methode}

\begin{astuce}
\textbf{Correspondance avec le dictionnaire classique :}

\begin{itemize}
    \item $\bar{\mathbf{b}} = B^{-1}\mathbf{b}$ : valeurs des variables de base (membre droit du dictionnaire)
    \item $\bar{A}_N = B^{-1}A_N$ : coefficients des variables non-basiques dans les equations
    \item $\bar{z} = \mathbf{c}_B^T B^{-1}\mathbf{b}$ : valeur courante de l'objectif
    \item $\bar{\mathbf{c}}_N = \mathbf{c}_N^T - \mathbf{c}_B^T B^{-1}A_N$ : coefficients reduits (ligne $z$ du dictionnaire)
\end{itemize}
\end{astuce}

\newpage
\subsection{Remarque sur l'Inverse de la Matrice de Base}

\begin{notecours}
\textbf{Ne pas calculer $B^{-1}$ explicitement !}

On pourrait penser que le simplexe revise necesssite de calculer l'inverse $B^{-1}$ explicitement. C'est une idee fausse pour plusieurs raisons :

\begin{itemize}
    \item Calculer $B^{-1}$ est \textcolor{importantrouge}{\textbf{plus long}} que calculer la factorisation triangulaire de $B$
    \item Les calculs intermediaires pour $B^{-1}$ introduisent des \textcolor{importantrouge}{\textbf{erreurs numeriques}} supplementaires
    \item Meme des matrices tres creuses (sparse) ont souvent des \textcolor{importantrouge}{\textbf{inverses denses}}
\end{itemize}

\textbf{En pratique :} On utilise une \textcolor{defvert}{\textbf{factorisation triangulaire}} de $B$ (elimination gaussienne) :
$$U = L_n P_n L_{n-1} P_{n-1} \cdots L_1 P_1 B$$

ou $P_i$ sont des matrices de permutation et $L_i$ des matrices triangulaires inferieures eta.

Pour resoudre $B\mathbf{d} = \mathbf{a}$, on resout successivement deux systemes triangulaires, ce qui est rapide.
\end{notecours}

\newpage
%=============================================
\section{Premier Exemple Detaille}
%=============================================

\needspace{10\baselineskip}
\subsection{Enonce et Mise en Forme}

\begin{exemple}
\textbf{Probleme de l'atelier de menuiserie :}

Un atelier fabrique quatre types de meubles : biblioth\`eques, bureaux, chaises et cadres de lit.

\begin{center}
\begin{tabular}{l|c|c|c|c}
\toprule
\textbf{Ressource} & Biblioth. & Bureau & Chaise & Cadre lit \\
\midrule
Heures de travail & 3 & 2 & 1 & 2 \\
Unites de metal & 1 & 1 & 1 & 1 \\
Unites de bois & 4 & 3 & 3 & 4 \\
\midrule
\textbf{Profit (\$)} & 19 & 13 & 12 & 17 \\
\bottomrule
\end{tabular}
\end{center}

\textbf{Disponibilites :} 225 heures de travail, 117 unites de metal, 420 unites de bois par jour.

\textbf{Formulation PL :}
\begin{align*}
\text{maximiser} \quad & z = 19x_1 + 13x_2 + 12x_3 + 17x_4\\
\text{sujet a} \quad & 3x_1 + 2x_2 + x_3 + 2x_4 \leq 225 \quad \text{(travail)}\\
& x_1 + x_2 + x_3 + x_4 \leq 117 \quad \text{(metal)}\\
& 4x_1 + 3x_2 + 3x_3 + 4x_4 \leq 420 \quad \text{(bois)}\\
& x_1, x_2, x_3, x_4 \geq 0
\end{align*}

\textbf{Variables d'ecart :} $x_5$ (travail), $x_6$ (metal), $x_7$ (bois)
\end{exemple}

\subsection{Donnees Matricielles}

\begin{exemple}
$$A = \begin{pmatrix} 3 & 2 & 1 & 2 & 1 & 0 & 0 \\ 1 & 1 & 1 & 1 & 0 & 1 & 0 \\ 4 & 3 & 3 & 4 & 0 & 0 & 1 \end{pmatrix}, \quad \mathbf{b} = \begin{pmatrix} 225 \\ 117 \\ 420 \end{pmatrix}, \quad \mathbf{c} = \begin{pmatrix} 19 \\ 13 \\ 12 \\ 17 \\ 0 \\ 0 \\ 0 \end{pmatrix}$$

\textbf{Solution de base initiale proposee :} $x_1 = 54, x_3 = 63$, autres variables nulles.

\textbf{Variables de base :} $x_1, x_3, x_7$ (et les autres hors-base). Donc :

$$B = \begin{pmatrix} 3 & 1 & 0 \\ 1 & 1 & 0 \\ 4 & 3 & 1 \end{pmatrix}, \quad \mathbf{x}_B = \begin{pmatrix} x_1 \\ x_3 \\ x_7 \end{pmatrix}, \quad \mathbf{c}_B = \begin{pmatrix} 19 \\ 12 \\ 0 \end{pmatrix}$$

Verification : $\bar{\mathbf{x}}_B = B^{-1}\mathbf{b} = (54, 63, 15)^T$ (tous $\geq 0$ : solution realisable \checkmark)
\end{exemple}

\newpage
%=============================================
\section{Une Iteration du Simplexe Revise}
%=============================================

\needspace{10\baselineskip}
\subsection{Structure d'une Iteration}

\begin{important}
\textbf{Les 5 etapes d'une iteration du simplexe revise :}

\textcolor{defvert}{\textbf{Test d'optimalite et variable entrante :}}

\begin{enumerate}
    \item[\circled{1}] \textbf{Resoudre} le systeme $\mathbf{y}^T B = \mathbf{c}_B^T$ pour trouver $\mathbf{y}^T = \mathbf{c}_B^T B^{-1}$

    \item[\circled{2}] \textbf{Calculer} $\mathbf{c}_N^T - \mathbf{y}^T A_N$ :
    \begin{itemize}
        \item Si tous les composantes sont $\leq 0$ : \textcolor{defvert}{\textbf{STOP -- solution optimale}}
        \item Sinon : choisir une colonne entrante $\mathbf{a}_j$ telle que $\mathbf{y}^T \mathbf{a}_j < c_j$
    \end{itemize}
\end{enumerate}

\textcolor{importantrouge}{\textbf{Test de borne et variable sortante :}}

\begin{enumerate}
    \item[\circled{3}] \textbf{Resoudre} le systeme $B\mathbf{d} = \mathbf{a}_j$ pour trouver $\mathbf{d} = B^{-1}\mathbf{a}_j$

    \item[\circled{4}] \textbf{Trouver} le plus grand $t$ tel que $\bar{\mathbf{x}}_B - t\mathbf{d} \geq \mathbf{0}$ :
    \begin{itemize}
        \item Si aucun $t$ n'existe (tous $d_i \leq 0$) : \textcolor{importantrouge}{\textbf{STOP -- probleme non borne}}
        \item Sinon : la variable sortante est celle qui impose la borne la plus stringente
    \end{itemize}
\end{enumerate}

\textcolor{algoviolet}{\textbf{Mise a jour de la base :}}

\begin{enumerate}
    \item[\circled{5}] \textbf{Mettre a jour} :
    \begin{itemize}
        \item Valeur de la variable entrante : $t$
        \item Valeurs des variables de base : $\bar{\mathbf{x}}_B \leftarrow \bar{\mathbf{x}}_B - t\mathbf{d}$
        \item Remplacer la \textcolor{importantrouge}{\textbf{colonne sortante}} de $B$ par la \textcolor{defvert}{\textbf{colonne entrante}}
        \item Mettre a jour le \textcolor{algoviolet}{\textbf{basis heading}}
    \end{itemize}
\end{enumerate}
\end{important}

\begin{definition}
\textbf{Basis Heading (En-tete de base) :}

La liste ordonnee des variables de base qui specifie l'ordre effectif des $m$ colonnes de $B$ est appelee le \textcolor{algoviolet}{\textbf{basis heading}}.

Par commodite, a chaque mise a jour, on remplace la variable sortante par la variable entrante dans le basis heading.
\end{definition}

\newpage
\subsection{Etape 1 et 2 : Variable Entrante}

\begin{methode}
\textbf{Calcul de $\mathbf{y}^T$ et test d'optimalite :}

\textbf{Etape \circled{1}} : Resoudre $\mathbf{y}^T B = \mathbf{c}_B^T$ (systeme triangulaire par substitution inverse).

\textbf{Etape \circled{2}} : Pour chaque variable non-basique $x_j$ (avec colonne $\mathbf{a}_j$), calculer le \textcolor{algoviolet}{\textbf{cout reduit}} :
$$\bar{c}_j = c_j - \mathbf{y}^T \mathbf{a}_j$$

\begin{itemize}
    \item Si $\bar{c}_j \leq 0$ pour tout $j \notin B$ : solution optimale
    \item Sinon : la \textcolor{defvert}{\textbf{variable entrante}} est tout $x_j$ avec $\bar{c}_j > 0$ (typiquement, on choisit celui avec le plus grand $\bar{c}_j$)
\end{itemize}

\textbf{Interpretation duale :} Le vecteur $\mathbf{y}^T = \mathbf{c}_B^T B^{-1}$ n'est autre que la \textcolor{importantrouge}{\textbf{solution duale courante}}.
\end{methode}

\begin{exemple}
\textbf{Application a notre exemple (base $B = \{x_1, x_3, x_7\}$) :}

\textbf{Etape \circled{1} :} Resoudre $[y_1 \; y_2 \; y_3] \begin{pmatrix} 3 & 1 & 0 \\ 1 & 1 & 0 \\ 4 & 3 & 1 \end{pmatrix} = [19 \; 12 \; 0]$

Par substitution : $y_3 = 0$, puis $y_1 + y_2 = 12$ et $3y_1 + y_2 = 19$, ce qui donne :

$$\mathbf{y}^T = \left[\frac{7}{2} \quad \frac{17}{2} \quad 0\right]$$

\textbf{Etape \circled{2} :} Pour les variables hors-base $x_2, x_4, x_5, x_6$ :

$$\mathbf{c}_N^T - \mathbf{y}^T A_N = [13 \; 17 \; 0 \; 0] - \left[\frac{7}{2} \; \frac{17}{2} \; 0\right] \begin{pmatrix} 2 & 2 & 1 & 0 \\ 1 & 1 & 0 & 1 \\ 3 & 4 & 0 & 0 \end{pmatrix} = \left[-\frac{5}{2} \quad \frac{3}{2} \quad -\frac{7}{2} \quad -\frac{17}{2}\right]$$

Seul $x_4$ a un cout reduit positif : $\bar{c}_4 = \frac{3}{2} > 0$. \textcolor{defvert}{\textbf{Variable entrante : $x_4$.}}
\end{exemple}

\newpage
\subsection{Etape 3 et 4 : Variable Sortante}

\begin{methode}
\textbf{Calcul de la direction de descente $\mathbf{d}$ :}

\textbf{Etape \circled{3} :} Resoudre $B\mathbf{d} = \mathbf{a}_j$ (colonne entrante de $A$).

Le vecteur $\mathbf{d} = B^{-1}\mathbf{a}_j$ donne les changements des variables de base quand la variable entrante augmente d'une unite.

\textbf{Etape \circled{4} :} La variable sortante est celle qui impose la borne la plus stringente sur $t$ :
$$t^* = \min_{i : d_i > 0} \frac{\bar{x}_{B_i}}{d_i}$$

La variable sortante correspond au minimum de ce rapport (regle du rapport minimum).
\end{methode}

\begin{exemple}
\textbf{Application (variable entrante $x_4$, colonne $\mathbf{a}_4 = (2, 1, 4)^T$) :}

\textbf{Etape \circled{3} :} Resoudre $B\mathbf{d} = \mathbf{a}_4$ :
$$\begin{pmatrix} 3 & 1 & 0 \\ 1 & 1 & 0 \\ 4 & 3 & 1 \end{pmatrix} \begin{pmatrix} d_1 \\ d_2 \\ d_3 \end{pmatrix} = \begin{pmatrix} 2 \\ 1 \\ 4 \end{pmatrix} \quad \Rightarrow \quad \mathbf{d} = \begin{pmatrix} 1/2 \\ 1/2 \\ 1/2 \end{pmatrix}$$

\textbf{Etape \circled{4} :} Ratios pour chaque variable de base :

\begin{center}
\begin{tabular}{c|c|c|c}
\toprule
Variable de base & Valeur $\bar{x}_{B_i}$ & $d_i$ & Ratio $\bar{x}_{B_i}/d_i$ \\
\midrule
$x_1 = 54$ & 54 & $1/2$ & $108$ \\
$x_3 = 63$ & 63 & $1/2$ & $126$ \\
$x_7 = 15$ & 15 & $1/2$ & \textcolor{importantrouge}{\textbf{30 (minimum)}} \\
\bottomrule
\end{tabular}
\end{center}

$t^* = 30$. \textcolor{importantrouge}{\textbf{Variable sortante : $x_7$.}}
\end{exemple}

\newpage
\subsection{Etape 5 : Mise a Jour de la Base}

\begin{methode}
\textbf{Mise a jour apres la premiere iteration :}

\textbf{Nouvelle base :} On remplace $x_7$ (sortant) par $x_4$ (entrant) :
\begin{itemize}
    \item Nouveau basis heading : $(x_1, x_3, x_4)$
    \item Valeur de la variable entrante : $x_4 = t^* = 30$
    \item Nouvelles valeurs des variables de base :
    $$\bar{\mathbf{x}}_B \leftarrow \bar{\mathbf{x}}_B - t^* \mathbf{d} = \begin{pmatrix} 54 \\ 63 \\ 15 \end{pmatrix} - 30 \begin{pmatrix} 1/2 \\ 1/2 \\ 1/2 \end{pmatrix} = \begin{pmatrix} 39 \\ 48 \\ 0 \end{pmatrix}$$
\end{itemize}

\textbf{Nouvelle matrice de base :} On remplace la colonne de $x_7$ dans $B$ par la colonne de $x_4$ :
$$B_{\text{nouveau}} = \begin{pmatrix} 3 & 1 & 2 \\ 1 & 1 & 1 \\ 4 & 3 & 4 \end{pmatrix}$$

\textbf{Nouvelle valeur de l'objectif :} $z = 19(39) + 12(48) + 17(30) = 741 + 576 + 510 = 1827$
\end{methode}

\begin{astuce}
\textbf{Remarque sur l'ordre des colonnes dans $B$ :}

L'ordre des colonnes de $B$ n'est important que si il correspond a l'ordre des composantes de $\mathbf{x}_B$.

Par exemple, le basis heading $(x_3, x_4, x_1)$ avec $\mathbf{x}_B = (x_3, x_4, x_1)^T = (48, 30, 39)^T$ et $B = \begin{pmatrix} 1 & 2 & 3 \\ 1 & 1 & 1 \\ 3 & 4 & 4 \end{pmatrix}$ est tout aussi valide.
\end{astuce}

\newpage
%=============================================
\section{Interpretation Economique}
%=============================================

\needspace{10\baselineskip}
\subsection{Signification des Variables Duales : Prix Fantomes}

\begin{notecours}
\textbf{Cadre economique :}

Dans un probleme de programmation lineaire typique :
\begin{itemize}
    \item $x_j$ mesure le \textcolor{titrebleu}{\textbf{niveau de production}} du produit $j$
    \item $b_i$ specifie la \textcolor{defvert}{\textbf{quantite disponible}} de la ressource $i$
    \item $a_{ij}$ : quantite de la ressource $i$ consommee par unite du produit $j$ (en unites de ressource $i$ par unite de produit $j$)
    \item $c_j$ : profit en dollars par unite du produit $j$
\end{itemize}

Les valeurs optimales duales $y_i^*$ ont alors une interpretation economique claire.
\end{notecours}

\begin{definition}
\textbf{Prix fantome (Shadow Price) :}

Pour $i = 1, 2, \ldots, m$, la variable duale optimale $y_i^*$ est appelee le \textcolor{importantrouge}{\textbf{prix fantome}} (shadow price) de la ressource $i$, ou encore sa \textcolor{defvert}{\textbf{valeur marginale}}.

Elle represente la \textcolor{algoviolet}{\textbf{variation du profit optimal}} par unite supplementaire de la ressource $i$ :
$$\frac{\partial z^*}{\partial b_i} = y_i^*$$

\textbf{Interpretation pratique :} Si la solution optimale est non degeneree, l'entreprise devrait etre prete a payer jusqu'a $y_i^*$ dollars \textbf{au-dela du prix de marche} pour obtenir une unite supplementaire de la ressource $i$.
\end{definition}

\subsection{Interpretation des Etapes du Simplexe Revise}

\begin{methode}
\textbf{Etape \circled{1} -- Prix fantomes temporaires :}

Resoudre $\mathbf{y}^T B = \mathbf{c}_B^T$ consiste a affecter des \textcolor{algoviolet}{\textbf{prix fantomes temporaires}} aux ressources tels que la \textbf{valeur totale des ressources consommees par chaque activite de base egale le profit} de cette activite.

\textbf{Pour notre exemple :} $y_1 = \frac{7}{2}$ (travail), $y_2 = \frac{17}{2}$ (metal), $y_3 = 0$ (bois)

\begin{itemize}
    \item Heure de travail : 3,50 \$/heure
    \item Unite de metal : 8,50 \$/unite
    \item Unite de bois : 0 \$/unite (ressource non saturee = valeur nulle)
\end{itemize}
\end{methode}

\begin{methode}
\textbf{Etape \circled{2} -- Evaluation des activites non-basiques (Pricing Out) :}

Calculer $\mathbf{y}^T A_N$ consiste a evaluer la valeur totale des ressources consommees par chaque activite non-basique. On compare cette valeur au profit qu'elle rapporte.

\begin{center}
\begin{tabular}{l|c|c|c}
\toprule
\textbf{Activite} & \textbf{Profit $c_j$} & \textbf{Cout ressources $\mathbf{y}^T \mathbf{a}_j$} & \textbf{Benefice net $\bar{c}_j$} \\
\midrule
Bureaux ($x_2$) & 13 & $\frac{7}{2}(2) + \frac{17}{2}(1) = \frac{31}{2}$ & $-\frac{5}{2} < 0$ \\
Cadres lit ($x_4$) & 17 & $\frac{7}{2}(2) + \frac{17}{2}(1) + 0(4) = \frac{31}{2}$ & $+\frac{3}{2} > 0$ \\
\bottomrule
\end{tabular}
\end{center}

La \textcolor{defvert}{\textbf{variable entrante}} est l'activite dont le \textbf{profit depasse le cout des ressources consommees}.
\end{methode}

\newpage
\begin{methode}
\textbf{Etape \circled{3} -- Le Mix :}

Resoudre $B\mathbf{d} = \mathbf{a}_j$ consiste a trouver la \textcolor{algoviolet}{\textbf{combinaison}} (mix) d'activites de base qui consomme les ressources au meme rythme que l'activite entrante.

\textbf{Pour notre exemple} (variable entrante : cadres de lit $x_4$) :

On cherche $d_1$ (biblio.), $d_3$ (chaises), $d_7$ (bois non utilise) tels que :
$$3d_1 + d_3 = 2, \quad d_1 + d_3 = 1, \quad 4d_1 + 3d_3 + d_7 = 4$$

Solution : $d_1 = \frac{1}{2}$, $d_3 = \frac{1}{2}$, $d_7 = \frac{1}{2}$.

Chaque cadre de lit sera substitue par : $\frac{1}{2}$ biblio. $+$ $\frac{1}{2}$ chaise $+$ $\frac{1}{2}$ unite de bois non utilise.
\end{methode}

\begin{definition}
\textbf{Substitution :}

L'operation consistant a remplacer l'activite entrante profitable par un mix approprie d'activites de base est appelee une \textcolor{defvert}{\textbf{substitution}}.

Une substitution augmente la valeur de l'objectif de :
$$c_j - \mathbf{y}^T \mathbf{a}_j = \bar{c}_j > 0$$

par unite d'activite entrante.

La valeur maximale $t$ que l'activite entrante peut prendre est determinee par la condition de ne pas consommer plus de ressources que disponible :
$$t^* = \max\{t : \bar{\mathbf{x}}_B - t\mathbf{d} \geq \mathbf{0}\}$$
\end{definition}

\begin{astuce}
\textbf{Couts Reduits (Reduced Costs) :}

Le \textcolor{algoviolet}{\textbf{cout reduit}} $\bar{c}_j = c_j - \mathbf{c}_B^T B^{-1} \mathbf{a}_j$ d'une variable non-basique $x_j$ represente la \textbf{variation nette de $z$} par unite d'augmentation de $x_j$ :

\begin{itemize}
    \item Augmenter $x_j$ d'une unite fait changer $\mathbf{x}_B$ de $-B^{-1}\mathbf{a}_j = -\mathbf{d}$
    \item La perte dans l'objectif due aux variables de base est $\mathbf{c}_B^T B^{-1}\mathbf{a}_j$
    \item Le gain direct est $c_j$
    \item Le \textbf{gain net} est $\bar{c}_j = c_j - \mathbf{c}_B^T B^{-1}\mathbf{a}_j$
\end{itemize}
\end{astuce}

\newpage
%=============================================
\section{Changement du Membre Droit et Prix Fantomes}
%=============================================

\needspace{10\baselineskip}
\subsection{Impact d'une Variation de $b_i$}

\begin{theoreme}
\textbf{Theoreme -- Sensibilite aux membres droits :}

Si le PL a au moins une solution de base optimale \textcolor{defvert}{\textbf{non degeneree}}, alors il existe $\varepsilon > 0$ tel que : si $|t_i| \leq \varepsilon$ pour tout $i = 1, \ldots, m$, le probleme perturbe

\begin{align*}
\text{maximiser} \quad & z = \sum_{j=1}^{n} c_j x_j\\
\text{sujet a} \quad & \sum_{j=1}^{n} a_{ij} x_j \leq b_i + t_i \quad \text{pour } i = 1, 2, \ldots, m\\
& x_j \geq 0
\end{align*}

admet une solution optimale de valeur :
$$z^*(t) = z^* + \sum_{i=1}^{m} y_i^* t_i$$

ou $z^*$ est la valeur optimale du probleme original et $y_1^*, y_2^*, \ldots, y_m^*$ est la solution duale optimale.
\end{theoreme}

\begin{important}
\textbf{Consequences pratiques :}

\begin{itemize}
    \item Le prix fantome $y_i^*$ donne la \textcolor{importantrouge}{\textbf{valeur exacte}} du changement de $z^*$ pour une petite variation de $b_i$
    \item Ce resultat est valable \textcolor{defvert}{\textbf{localement}} (pour des petites perturbations $|t_i| \leq \varepsilon$)
    \item Pour des perturbations plus grandes, la base optimale peut changer (on sort de la region de validite)
    \item En cas de \textcolor{algoviolet}{\textbf{degenerescence}}, plusieurs prix fantomes peuvent correspondre a la meme solution primale optimale
\end{itemize}
\end{important}

\subsection{Exemple : L'Application Forestiere}

\begin{exemple}
\textbf{Probleme du forestier :}

Un forestier dispose de 100 acres de bois de feuillus. Deux options :
\begin{itemize}
    \item \textbf{Option 1 :} Couper et laisser regenerer naturellement -- cout immediat 10\$/acre, retour 50\$/acre (profit net : 40\$/acre)
    \item \textbf{Option 2 :} Couper et replanter en pins -- cout immediat 50\$/acre, retour 120\$/acre (profit net : 70\$/acre)
\end{itemize}

Budget immediat disponible : 4 000\$.

\textbf{Formulation PL} (avec $x_1$ = acres option 1, $x_2$ = acres option 2) :
\begin{align*}
\text{maximiser} \quad & z = 40x_1 + 70x_2\\
\text{sujet a} \quad & x_1 + x_2 \leq 100 \quad \text{(superficie)}\\
& 10x_1 + 50x_2 \leq 4000 \quad \text{(budget)}\\
& x_1, x_2 \geq 0
\end{align*}

\textbf{Solution optimale :} $x_1^* = 25$ acres, $x_2^* = 75$ acres, $z^* = 25(40) + 75(70) = 1000 + 5250 = 6250$ \$.
\end{exemple}

\newpage
\begin{exemple}
\textbf{Calcul des prix fantomes :}

Dual optimal : $y_1^*$ (superficie) et $y_2^*$ (budget).

Par les ecarts complementaires (les deux contraintes sont saturees) :
\begin{align*}
y_1^* + y_2^* &= 40 \quad \text{(dual contrainte } x_1 > 0\text{)}\\
y_1^* + 5y_2^* &= 70 \quad \text{(dual contrainte } x_2 > 0\text{, coeff. 50/10 = 5)}
\end{align*}

\textbf{Attention :} La deuxieme contrainte duale est $y_1^* + \frac{50}{10}y_2^* \geq 70$, mais en divisant par les bonnes unites... Utilisons directement $y_1^*$ = prix par acre et $y_2^*$ = prix par dollar de budget :

$$y_1^* = 22.5, \quad y_2^* = \frac{17.5}{10} = 1.75 \text{ par dollar} \Rightarrow y_2^* = 0.175 \text{ par dollar}$$

\textbf{Questions sur la perturbation :}

\textbf{1. Emprunter 100\$ maintenant pour rembourser 180\$ plus tard ?}

Emprunter augmente le budget de 100\$ ($t_2 = +100$) et reduit le profit futur de 180\$. L'effet sur $z^*$ :
$$\Delta z = y_2^* \cdot 100 - 180 = 0.175 \times 100 \times 10 - 180$$

\textit{(Nota : $y_2^*$ est en profit par dollar de budget, le budget est en dollars)}

Gain brut : $y_2^* \times 100 = 17.5$ \$ de profit supplementaire. Cout : 180\$. 

\textcolor{importantrouge}{\textbf{Conclusion :}} $17.5 < 180$. Ne pas emprunter.

\textbf{2. Investir 100\$ maintenant pour recevoir 180\$ plus tard ?}

Investir reduit le budget de 100\$ ($t_2 = -100$) et rapporte 180\$ plus tard.

Perte : $y_2^* \times 100 = 17.5$ \$ de profit perdu. Gain : 180\$.

\textcolor{defvert}{\textbf{Conclusion :}} $180 > 17.5$. Investir est profitable.
\end{exemple}

\begin{astuce}
\textbf{Interpretation generale :}

\begin{itemize}
    \item Toute operation qui \textcolor{defvert}{\textbf{rapporte plus}} que $y_i^*$ par unite de ressource $i$ est \textbf{profitable}
    \item Toute operation qui \textcolor{importantrouge}{\textbf{coute plus}} que $y_i^*$ par unite de ressource $i$ est \textbf{non profitable}
    \item Les prix fantomes fournissent les \textcolor{algoviolet}{\textbf{prix d'equilibre}} qui maximisent le profit global
\end{itemize}
\end{astuce}

\newpage
%=============================================
\section{Resume et Comparaison des Methodes}
%=============================================

\begin{important}
\textbf{Simplexe Standard vs. Simplexe Revise :}

\begin{center}
\begin{tabular}{l|p{5cm}|p{5cm}}
\toprule
& \textbf{Simplexe Standard} & \textbf{Simplexe Revise} \\
\midrule
\textbf{Stockage} & Tableau complet $m \times (n+m)$ & Matrice de base $B$ ($m \times m$) \\
\textbf{Calculs par iteration} & Mise a jour du tableau entier & 2 systemes lineaires \\
\textbf{Donnees utilisees} & Tableau transforme & Donnees originales $A, b, c$ \\
\textbf{Avantage} & Simple a implementer & Rapide pour grands problemes \\
\textbf{Usage pratique} & Pedagogique & Tous les logiciels modernes \\
\bottomrule
\end{tabular}
\end{center}
\end{important}

\begin{astuce}
\textbf{Recapitulatif des formules cles :}

\begin{center}
\begin{tabular}{l|l}
\toprule
\textbf{Element} & \textbf{Formule matricielle} \\
\midrule
Valeurs variables de base & $\bar{\mathbf{x}}_B = B^{-1}\mathbf{b}$ \\
Solution duale courante & $\mathbf{y}^T = \mathbf{c}_B^T B^{-1}$ (resoudre $\mathbf{y}^T B = \mathbf{c}_B^T$) \\
Couts reduits & $\bar{c}_j = c_j - \mathbf{y}^T \mathbf{a}_j$ \\
Direction $\mathbf{d}$ (entrant $x_j$) & $\mathbf{d} = B^{-1}\mathbf{a}_j$ (resoudre $B\mathbf{d} = \mathbf{a}_j$) \\
Ratio test & $t^* = \min_{d_i > 0} \bar{x}_{B_i}/d_i$ \\
Valeur objectif & $z = \mathbf{c}_B^T B^{-1}\mathbf{b} = \mathbf{y}^T \mathbf{b}$ \\
Sensibilite $b$ & $z^*(b + t) \approx z^* + \mathbf{y}^{*T} \mathbf{t}$ \\
\bottomrule
\end{tabular}
\end{center}
\end{astuce}

\newpage
%=============================================
\section{Exercices}
%=============================================

\begin{exemple}
\textbf{Exercice 1 -- Formulation matricielle}

Pour le PL suivant :
\begin{align*}
\text{maximiser} \quad & z = 3x_1 + 2x_2 + 5x_3\\
\text{sujet a} \quad & x_1 + 2x_2 + x_3 \leq 6\\
& 2x_1 + x_2 + 3x_3 \leq 10\\
& x_1, x_2, x_3 \geq 0
\end{align*}

Avec les variables d'ecart $x_4, x_5$ et la base $B = \{x_1, x_3\}$ :

\begin{enumerate}
    \item Ecrire la matrice de base $B$, les vecteurs $\mathbf{c}_B$, $\mathbf{c}_N$, la matrice $A_N$
    \item Calculer les valeurs des variables de base $\bar{\mathbf{x}}_B = B^{-1}\mathbf{b}$
    \item Calculer la solution duale courante $\mathbf{y}^T = \mathbf{c}_B^T B^{-1}$
    \item Calculer les couts reduits des variables hors-base
    \item Conclure sur l'optimalite ou identifier la variable entrante
\end{enumerate}
\end{exemple}

\begin{exemple}
\textbf{Exercice 2 -- Iteration complete du simplexe revise}

En partant de la solution initiale (variables d'ecart en base) pour le PL :
\begin{align*}
\text{maximiser} \quad & z = 4x_1 + 3x_2\\
\text{sujet a} \quad & x_1 + x_2 \leq 4\\
& 2x_1 + x_2 \leq 6\\
& x_1, x_2 \geq 0
\end{align*}

Effectuer \textbf{deux iterations completes} du simplexe revise en suivant les 5 etapes. A chaque iteration, verifier la valeur de $z$.
\end{exemple}

\begin{exemple}
\textbf{Exercice 3 -- Prix fantomes et analyse de sensibilite}

On revient au probleme de l'atelier (Section 3.1). La solution optimale est $x_1^* = 54, x_3^* = 63, x_4^* = 30$ avec $z^* = 1827$ et les prix fantomes $y_1^* = \frac{7}{2}$, $y_2^* = \frac{17}{2}$, $y_3^* = 0$.

\begin{enumerate}
    \item Interpreter economiquement chaque prix fantome
    \item L'entreprise peut obtenir 5 heures de travail supplementaires au prix de 10\$/heure. Vaut-il mieux les acheter ?
    \item L'entreprise peut vendre 2 unites de metal a 6\$ l'unite. Doit-elle le faire ?
    \item Quel est l'impact sur $z^*$ si le bois disponible passe de 420 a 430 unites ?
\end{enumerate}
\end{exemple}

\begin{exemple}
\textbf{Exercice 4 -- Interpretation economique des couts reduits}

Expliquer pourquoi, dans la premiere iteration de notre exemple, le cout reduit de $x_2$ (bureaux) est negatif alors que le profit unitaire des bureaux est positif (\$13). Comment cela est-il coherent avec la theorie ?
\end{exemple}

\end{document}
